\documentclass{article}
\usepackage[utf8]{inputenc}
\usepackage[T1]{fontenc}
\usepackage{geometry}
\geometry{
    a4paper,
    left=2.5cm,
    right=2.5cm,
    top=2.5cm,
    bottom=2.5cm
}
\usepackage{graphicx}
\usepackage{float}
\usepackage{hyperref}
\usepackage[english]{babel}
\usepackage{amsmath}
\usepackage{amssymb}
\usepackage{array}
\usepackage{booktabs}
\usepackage{caption}
\usepackage{subcaption}
\usepackage{enumitem}
\usepackage{fancyhdr}
\usepackage{titlesec}
\usepackage{url}
\usepackage{textcomp}
\usepackage{xcolor}
\usepackage{listings}
\usepackage{multirow}
\usepackage{longtable}
\usepackage{tabularx}
\usepackage{times}

% Set up page style
\pagestyle{fancy}
\fancyhf{}
\rfoot{\thepage}
\renewcommand{\headrulewidth}{0pt}

% Title formatting
\titleformat{\chapter}[display]{\normalfont\bfseries}{\chaptertitle\thechapter}{0pt}{\Huge}
\titlespacing*{\chapter}{0pt}{50pt}{40pt}

% Document info
\title{Template Document}
\author{Academic Support Agent}
\date{\today}

\begin{document}
\maketitle

\przedmiot{Przetwarzanie Sygna w w Systemach Diagnostyki Medycznej}
\tytul{Laboratorium 6}
\podtytul{Analiza sygna u w dziedzinie czasowo-cz stotliwo ciowej z zastosowaniem
transformacji falkowych}
\kierunek{Informatyka Stosowana}
\autor{Przemys aw Elias}
\data{Krak w, 1 czerwca 2011}

\stronatytulowa{}

\section{Wprowadzenie}

Dokument przedstawia przyk adowe u ycie klasy \emph{sprawozdanie-agh} do pisania sprawozda  w systemie \LaTeX.

\section{Konfiguracja}

Do poprawnej kompilacji wymagane s  nast puj ce pliki:
\begin{enumerate}
\item \texttt{sprawozdanie-agh.cls}
\item \texttt{logo-agh.jpg}
\end{enumerate}

Pliki te mo esz skopiowa  do katalogu, w kt rym znajduje si  dokument \texttt{*.tex} (np.: taki jak ten przyk adowy). Wad  tego rozwi zania jest konieczno  ka dorazowego kopiowania tych e plik w w miejsce, gdzie znajduje si  dokument, kt ry chcemy skompilowa .\\
Zalecam zatem skopiowanie plik w \texttt{sprawozdanie-agh.cls} oraz \texttt{agh-logo.jpg} do zbioru klas \LaTeX, dzi ki czemu nie b dzie konieczno ci ich ka dorazowego kopiowania. Je li korzystasz z systemu MiKTeX pliki te powiniene  umie ci  w katalogu:
\begin{lstlisting}
C:\Program Files\MiKTeX 2.9\tex\latex\AGH
\end{lstlisting}
 cie ka ta mo e si  r ni  dla Twojej instalacji. Tak jak poda em w przyk adowej  cie ce, warto stworzy  osobny katalog \texttt{AGH} i dopiero w nim umie ci  omawiane dwa pliki.
\end{document}